\begin{abstract}
Large language model agents are increasingly deployed to browse the web
autonomously on behalf of users, acting on tasks such as job applications,
online purchases, and form submissions. To complete such tasks without
constant supervision, these agents are granted access to the user's
personally identifiable information (PII), including names, addresses,
payment details, and government identifiers. This access creates a
confused-deputy attack surface: adversarial web content can coerce an
agent into forwarding sensitive data to an attacker-controlled endpoint.
Existing benchmarks evaluate agent susceptibility to dark patterns and
prompt injection, measuring task-trajectory divergence rather than
data exfiltration, and none incorporate multi-turn conversational attacks
or a structured PII profile. We introduce \textsc{Scammer4U},
a benchmark of 91 adversarial environments plus 10 benign-twin baselines,
spanning 16 site categories and 8 attack vectors, organized on an
8-axis factorial taxonomy that supports controlled paired-sibling
ablation studies. We evaluate four frontier model families
(GPT-5 mini, Claude Haiku 4.5, Gemini 3 Flash, and Llama 4 Scout) on
four metrics: PII Leakage Rate (PLR), Attack Success Rate (ASR),
Task Completion Rate (TCR), and Detection Rate (DR); the cross-family
LLM-as-judge that operationalises DR has completed on three of the
four panels at submission, so the detection-conditioned result below
is pooled over those three.
Within the pre-registered C3 pre-submission reflection condition,
restricted to sessions that reach the attack surface, agents whose
reasoning trace an independent LLM-as-judge confirms has flagged the
environment as suspicious still submit critical PII in 38.0\% of
sessions, versus 69.5\% when no such recognition is verbalised: a
31.5-percentage-point detection--action gap pooled across the three
model families on which the cross-family judge has been completed
(Claude Haiku 4.5, Gemini 3 Flash, Llama 4 Scout; $n = 295$
DR$=$1 sessions), statistically robust under the pre-registered
analysis. Our results expose fundamental safety gaps
in current autonomous web agents for PII-sensitive tasks and motivate
output-level defenses that do not rely on the agent's ability to
recognise an attack before acting.
\end{abstract}
